% Folder 91, batch 04 — archivists' pages 61–80.
% Pass: Opus 5.5 (claude-opus-5-5), 2026-10-03 — first pass, unchecked against the pages by a human.
%
% Source: archives/batches/91/batch-04.pdf (cover sheet excluded), read from
% the 700-dpi tiles in archives/tiles/91/p61 … p80. Fast draft hand: the
% formulas mostly carry, the connective prose often does not. Page 65 is a
% double sheet (two columns). Page 80 is a cover in his hand naming the next
% run; it gets no \page{}, its words become the closing heading.
\documentclass[11pt,a4paper]{article}
\input{../preamble/grothendieck}

\folder{91}
\batch{4}
\pages{61}{80}
\foldertitle{Autour de Néron / Greenberg-Néron. Foncteurs Hom (méthodes non-projectives) : notes manuscrites (s.d.), lettre (1967).}
\dating{[à partir de 1964-vers 1970]}
\watermark{Édition de démonstration}

\begin{document}

\page{61}
\note{la page continue une liste d'exemples commencée avant ce lot (exemples I et II, auxquels elle renvoie). Notation du lot : le foncteur que Grothendieck écrit d'un W barré et souligné est transcrit $\mathbb{W}$ ; ses préschémas en lettres d'écriture $\mathfrak{X}$, $\mathfrak{Y}$, $\mathfrak{U}$ ; l'adjoint qu'il note d'un exposant $a$ est $\mathbb{W}^{a}$.}

\struck{\ill{}} et que l'homomorphisme \uncertain{naturel}
$\Gamma(W_T/T) \to \Gamma(W_{T'}/T')$ pour $T' \to T$ \struck{défini} soit
\uncertain{local} sur les \ill{} faisceaux. Les données \add{$\mathbb{W}$}
\ill{} ainsi un foncteur \uncertain{covariant} sur les $S$-préschémas, et un
isomorphisme fonctoriel $|\mathbb{W}(T)| \simeq |T|$ (\ill{}
\uncertain{sous-jacents}). \ill{} \struck{\ill{} $\operatorname{Hom}_{S}(T,\Lambda) =$ \ill{}}
\note{une ligne biffée et une grande boucle de rature couvrent ici deux lignes, où se lisent seulement « $\operatorname{Hom}(T,\Lambda)=$ » et « \ill{} ait un \ill{} ».}

\struck{$\Lambda$, puis \ill{}} \ill{} \uncertain{suffisant} \ill{}.

\textbf{Exemple III} Soit $S' = \operatorname{Spec}(\mathcal{A})$, où
$\mathcal{A}$ est un faisceau \uncertain{loc.\ libre} d'algèbres \ill{} sur
$\mathcal{O}_S$. \ill{} \uncertain{plus} \ill{} $S' \to S$ \ill{}
\uncertain{morphisme} \uncertain{universel}. \ill{}
$\Lambda = \struck{\ill{}} \mathbf{V}(\check{\mathcal{A}})$. On est dans
\ill{} \uncertain{rien} \ill{} de l'exemple I et de l'exemple II.
\marginal{$S' \to S$}

\textbf{Exemple III$'$} Comme dans l'exemple II, mais \ill{} \ill{}
\ill{} $\mathbb{W} \to S[t]$ d'\uncertain{anneaux}, tel que pour tout
$T/S$, l'\uncertain{application} \uncertain{correspondante}
$\mathcal{O}_{\mathbb{W}(T)} \to \mathcal{O}_T$ soit surjective, de noyau
\uncertain{nilpotent}. Soit $\mathfrak{X}$ un préschéma sur
$\mathbb{W}(S)$. Considérons \ill{} le foncteur \uncertain{contravariant}
sur les $S$-préschémas $T$
\[
  (*) \qquad \operatorname{Hom}_{\mathbb{W}(S)}(\mathbb{W}(T), \mathfrak{X}) .
\]
\marginal{\uncertain{Cas} \uncertain{particulier} de l'exemple II :
\uncertain{foncteurs} de Greenberg}

\page{62}
On \struck{\ill{}} le \uncertain{désignera} \ill{} s'il est représentable,
i.e.\ \ill{} \ill{} le \uncertain{schéma} sur $T$ \ill{} foncteur
\uncertain{adjoint} (au sens de Kan) du foncteur $\mathbb{W}$ de
$S$-préschémas en $T$-préschémas. Une solution des pbs
\uncertain{universels} posés par le foncteur $\mathbb{W}$ \ill{} $T$
\ill{} \ill{} est un couple ($X$ $S$-préschéma) $\mathbb{W}(X)
\xrightarrow{\alpha} \mathfrak{X}$ tel que pour tout $T$, l'application
naturelle
\[
  \operatorname{Hom}_{S}(T, X) \xrightarrow{\ \sim\ }
  \operatorname{Hom}_{\mathbb{W}(S)}(\mathbb{W}(T), \mathfrak{X})
\]
\struck{soit \ill{}} définie par $u \mapsto \bigl(\mathbb{W}(T)
\xrightarrow{\mathbb{W}(u)} \mathbb{W}(X) \xrightarrow{\alpha}
\mathfrak{X}\bigr)$, soit \emph{bijective}.

\ill{} (*) \ill{} $T \to$ \ill{} « \uncertain{foncteur} \uncertain{local} »,
il suffit de \uncertain{vérifier} la \ill{} \ill{} \ill{}, $T$
\uncertain{affine} \ill{} \uncertain{affine} \uncertain{sur} $S$,
\uncertain{bijectivité} \ill{}, p.\ ex.

Désignons par $\mathbb{W}^{a}$ le foncteur \uncertain{adjoint} de
$\mathbb{W}$ (défini sur une \ill{} \ill{} sous-catégorie de la catégorie
des $T$-préschémas)
\[
  \operatorname{Hom}_{S^{a}}(T^{a}, \mathfrak{X}) \simeq
  \operatorname{Hom}_{S}(T, \mathfrak{X})
\]
\note{la formule finale est en bas de page, coupée au bord ; ses indices et exposants ($S^{a}$, $T^{a}$) sont incertains.}

\page{63}
(i) Si $\mathbb{W}^{a}(\mathfrak{X})$ et $\mathbb{W}^{a}(\mathfrak{Y})$ sont
définis, il en est de même de $\mathbb{W}^{a}(\mathfrak{X}
\times_{\mathbb{W}(S)} \mathfrak{Y})$, et on a
\[
  \mathbb{W}^{a}(\mathfrak{X} \times_{\mathbb{W}(S)} \mathfrak{Y}) \simeq
  \mathbb{W}^{a}(\mathfrak{X}) \times_{S} \mathbb{W}^{a}(\mathfrak{Y})
\]
et \ill{} plus généralement \struck{pour} la formule analogue pour les limites
projectives finies [\uncertain{donc} ($\alpha$) l'objet \uncertain{final}
($\beta$) les produits fibrés, d'où comme cas particulier ($\gamma$) les
noyaux].

(ii)
\[
  \mathbb{W}^{a}(\mathfrak{X}) \times_{S} T \simeq
  \mathbb{W}_{T}^{a}(\mathfrak{X} \times_{\mathbb{W}(S)} \mathbb{W}(T))
\]
[l'existence du premier membre \uncertain{implique} celle du second ; la
définition de $\mathbb{W}_T$ est claire \ill{}]

De (i) résulte que le foncteur $\mathbb{W}^{a}$ transforme un préschéma en
groupes, \uncertain{en} anneaux etc.\ en un préschéma du même type.

\page{64}
\uncertain{Supposons} que $\mathbb{W}^{a}(\mathbb{W}(S)[t])$ existe. C'est
alors, de façon naturelle, un préschéma en anneaux sur $S$, que nous
\uncertain{noterons} $\mathbb{W}$, appelé \emph{préschéma d'anneaux
\uncertain{générateur} du foncteur $\mathbb{W}$}.

[Dans l'exemple II, c'est le schéma d'anneaux $W$ de \uncertain{départ}.
Dans le cas de l'exemple I, \struck{\ill{}} si \struck{\ill{}} on suppose
$S'$ \add{fini et} loc.\ libre sur $S$, c'est le \struck{\ill{}} schéma en
anneaux $\mathbf{V}(\check{\mathcal{A}})$, où $\mathcal{A}$ est le faisceau
d'algèbres sur $S$ \uncertain{associé} à $S'$].

Soit alors $\mathfrak{X}$ \uncertain{sous-schéma} \uncertain{fermé} de
$\mathbb{W}(S)[t_1, \ldots, t_r]$ défini par \struck{\ill{}} le faisceau
d'idéaux dans $\mathcal{O}_{\mathbb{W}(S)}[t_1, \ldots, t_r]$ engendré par
des sections $F_i$ ($i = 1, \ldots, s$), donc $\mathfrak{X}$ est le
\uncertain{produit} fibré, sur $\mathcal{Y}^{s}$, de $\mathcal{Y}^{r}$ et
$\mathcal{Y}^{0}$ [où $\mathcal{Y}^{*} = \mathbb{W}[S][t]$].

\page{65}
\note{feuille double, en deux colonnes ; la colonne de gauche d'abord.}
Alors $\mathbb{W}^{a}(\mathfrak{X})$ existe, et \ill{} \ill{} \ill{} comme
produit fibré. De façon \uncertain{précise}, les $F_j$ (\add{$1 \le j \le s$})
(\uncertain{définissent} des polynômes \ill{} les $t_i$ ($1 \le i \le r$), à
coefficients dans les $\Gamma(\Lambda/S)$, d'où \struck{\ill{}} un
$S$-morphisme $\Lambda^{r} \to \Lambda^{s}$, et $\mathbb{W}^{a}(\mathfrak{X})$
est \uncertain{isomorphe} à l'image \uncertain{inverse} de la section
\uncertain{nulle} de $\Lambda^{s}$.

\textbf{Exemples} \struck{\ill{}} $(\mathbb{G}_{a\,\mathbb{W}(S)})^{n}$,
$\mathrm{Gl}(n)_{\mathbb{W}(S)}$, etc.

Soit $\mathbb{W}'(T) \xrightarrow{\alpha(T)} \mathbb{W}(T)$ un morphisme
fonctoriel de préschémas, définissant en particulier
\[
  \mathbb{W}'(S) \to \mathbb{W}(S) .
\]
Soit de plus $\mathfrak{X}$ sur $\mathbb{W}(S)$, \uncertain{posons}
\[
  \mathfrak{X}' = \mathfrak{X} \times_{\mathbb{W}(S)} \mathbb{W}'(S) .
\]
On a donc, pour tout $T$, une \uncertain{application}
\add{\struck{\ill{}}}
\[
  (*) \qquad \operatorname{Hom}_{\mathbb{W}(S)}(\mathbb{W}(T), \mathfrak{X})
  \to \operatorname{Hom}_{\mathbb{W}'(S)}(\mathbb{W}(T)
  \times_{\mathbb{W}(S)} \mathbb{W}'(S), \mathfrak{X}')
\]
(l'image réciproque d'un morphisme par $\mathbb{W}'(S) \to \mathbb{W}(S)$]
d'autre part, on a
\[
  \mathbb{W}'(T) \to \mathbb{W}(T) \times_{\mathbb{W}(S)} \mathbb{W}'(S)
\]
\struck{\ill{}} grâce aux morphismes $\mathbb{W}'(T) \to
\struck{\ill{}} \mathbb{W}(T)$ et $\mathbb{W}'(T) \to \mathbb{W}'(S)$,
d'où
\[
  (**) \qquad \operatorname{Hom}_{\mathbb{W}'(S)}(\mathbb{W}(T)
  \times_{\mathbb{W}(S)} \mathbb{W}'(S), \mathfrak{X}') \to
  \operatorname{Hom}_{\mathbb{W}'(S)}(\mathbb{W}'(T), \mathfrak{X}')
\]
d'où en composant (*) et (**) un \add{morphisme} \struck{\ill{}},
\uncertain{évidemment} \uncertain{fonctoriel} en $T$
\[
  \operatorname{Hom}_{\mathbb{W}(S)}(\mathbb{W}(T), \mathfrak{X})
  \longrightarrow
  \operatorname{Hom}_{\mathbb{W}'(S)}(\mathbb{W}'(T), \mathfrak{X}') .
\]
Par suite, si $\mathbb{W}^{a}(\mathfrak{X})$ et
$\mathbb{W}'^{a}(\mathfrak{X}')$ existent, on a un morphisme canonique
\[
  \mathbb{W}^{a}(\mathfrak{X}) \xrightarrow{\alpha^{*}(\mathfrak{X})}
  \mathbb{W}'^{a}(\mathfrak{X} \times_{\mathbb{W}(S)} \mathbb{W}'(S)) .
\]

\page{66}
De cette façon, pour $\mathbb{W}$ fixé, \struck{\ill{}} \uncertain{lorsque}
$\mathfrak{X}$ et $\mathbb{W}'$ varient,
\[
  \mathbb{W}'^{a}(\mathfrak{X} \times_{\mathbb{W}(S)} \mathbb{W}'(S))
\]
devient un \emph{bifoncteur} en $\mathfrak{X}$, $\mathbb{W}'$,
\emph{covariant} en $\mathfrak{X}$, \emph{contravariant} en $\mathbb{W}'$.

\textbf{(2) Critère d'existence : cas affine}

Supposons que $\mathbb{W}$ transforme schémas affines (sur \ill{} affine
de $S$) en schémas affines. Donc, \uncertain{au-dessus} d'un \ill{} affine
de $S$, d'anneau $A$, $\mathbb{W}$ est défini par un foncteur
\uncertain{covariant}, \uncertain{noté} aussi $\mathbb{W}$, transformant
\add{$B$} $A$-algèbres en anneaux $\mathbb{W}(B)$. [\struck{\ill{}}
L'hypothèse de \struck{\uncertain{localité}} \add{\uncertain{caractère} local}
de \ill{} foncteur, \ill{} \ill{} \uncertain{explicitation} \ill{}]. Le
\ill{} \ill{} \ill{} \ill{}
\[
  \mathbb{W}(B) = \operatorname{Hom}_{A\text{-}\mathrm{alg}}(\Lambda, B)
\]
où $\Lambda$ est \struck{l'algèbre affine} \add{l'anneau} d'un schéma en
anneaux \struck{affine} sur $S$, affine sur $S$.

\marginal{\uncertain{Ici} il suffit que $\mathbb{W}$ \uncertain{transforme}
\ill{} affines \struck{\ill{}} \ill{}. Foncteur $\mathbb{W}$
\emph{\uncertain{spécial}} \ill{} \struck{\ill{}} \ill{} \uncertain{schémas}
affines sur $S$}

\page{67}
[C'est \ill{} le cas en particulier dans le cas de l'exemple I,
\struck{\ill{}} \add{lorsque $S'$ est} fini et loc.\ libre sur $S$, et
l'exemple II, \add{lorsque} $W$ \struck{\ill{}} affine sur $S$].

\textbf{Proposition} Sous ces conditions, si \struck{\ill{}}
$\mathfrak{X}$ est un schéma \emph{affine} sur $\mathbb{W}(S)$,
$\mathbb{W}^{a}(\mathfrak{X})$ existe et est un schéma \emph{affine} sur
$S$.

\textit{Démonstration.} \struck{Soit $\mathfrak{X}$} \add{\uncertain{on peut
supposer} $S$ affine}. \uncertain{Soit} $\mathfrak{X} =
\operatorname{Spec}(\mathcal{C})$ où $\mathcal{C}$ est une algèbre de type
fini \uncertain{de} \uncertain{présentation} sur $\mathbb{W}(A)$,
\uncertain{engendrée} \ill{} \ill{}, il résulte des résultats \ill{} \ill{}
que $\mathbb{W}^{a}(\mathfrak{X})$ existe et est \ill{} \ill{}
\uncertain{sous-préschéma} \uncertain{fermé} de $\mathbb{W}^{r}$ [défini
par des « équations \uncertain{polynomiales} » \uncertain{à coeff.\ dans}
$\Gamma(\mathbb{W}/S)$]. \add{\ill{}} \ill{} $\mathbb{W}$ \ill{}
\uncertain{affine}, il en est de même de \struck{\ill{}} $\mathbb{W}^{r}$
\ill{} \ill{} de $\mathbb{W}^{a}(\mathfrak{X})$ (\uncertain{cqfd}). \ill{}
\uncertain{obtenu} \ill{}

\textbf{Cor.} Si \struck{\ill{}} \add{$S$-\ill{}, $S$ affine}. Si
$\mathfrak{X}$ est \ill{} \ill{} \ill{} \uncertain{générateurs}
\ill{} sur celui de $\mathbb{W}(S)$, et celui de $\mathbb{W}^{a}$ \ill{}
\uncertain{engendré} sur celui de $S$, alors

\page{68}
l'anneau de $\mathbb{W}^{a}(\mathfrak{X})$ est engendré par \ill{}
\ill{} générateurs sur celui de $S$.

Dans le cas général \add{($S$ affine)}, on a $\mathcal{C} =
\varinjlim \mathcal{C}_i$, où les $\mathcal{C}_i$ sont des
$\mathbb{W}(A)$-algèbres \add{de type fini} \uncertain{dans} $\mathcal{C}$.
On a donc, \uncertain{pour} $T$ affine \ill{}
\[
\begin{aligned}
  \operatorname{Hom}_{\mathbb{W}(S)}(\mathbb{W}(T), \mathfrak{X})
  &= \operatorname{Hom}_{\mathbb{W}(A)\text{-}\mathrm{alg}}(\mathcal{C},
     \mathbb{W}(B)) \\
  &\simeq \varprojlim \operatorname{Hom}_{\mathbb{W}(A)\text{-}\mathrm{alg}}
     (\mathcal{C}_i, \mathbb{W}(B)) \\
  &\simeq \varprojlim \operatorname{Hom}_{A\text{-}\mathrm{alg}}
     (\mathbb{W}^{a}(\mathcal{C}_i), B) \\
  &\simeq \operatorname{Hom}_{A\text{-}\mathrm{alg}}
     (\varinjlim \mathbb{W}^{a}(\mathcal{C}_i), B)
\end{aligned}
\]
\note{dans la première ligne, sous l'argument $\mathcal{C}$, un premier argument est biffé ($\mathbb{W}(B)$ ?).}

Soit donc
\[
  \varinjlim \mathbb{W}^{a}(\mathcal{C}_i) = \mathcal{D},
\]
\struck{\ill{}} on obtient un isomorphisme fonctoriel
\[
  \operatorname{Hom}_{\mathbb{W}(S)}(\mathbb{W}(T), \mathfrak{X}) \simeq
  \operatorname{Hom}_{S}(T, \operatorname{Spec}(\mathcal{D}))
\]
d'où le résultat.

\struck{\textbf{Corollaire} Sous les conditions précédentes, si
$\mathfrak{X} \xrightarrow{f} \mathfrak{Y}$ est un morphisme
\struck{\ill{}} de $S$-préschémas affines \ill{} \ill{} \ill{}
(\uncertain{c.-à-d.} \ill{} \ill{}) \ill{} \ill{}}

\struck{\textbf{Corollaire 2} Soit $f\colon \mathfrak{X} \to \mathfrak{Y}$
\ill{} $\mathbb{W}(S)$-préschémas, \ill{} \ill{} que
$\mathbb{W}^{a}(\mathfrak{X})$ et $\mathbb{W}^{a}(\mathfrak{Y})$
existent, alors $\mathbb{W}^{a}(f)\colon$}
\struck{$\mathbb{W}^{a}(\mathfrak{X}) \to \mathbb{W}^{a}(\mathfrak{Y})$
\ill{}}
\note{ces deux énoncés sont barrés de longues diagonales ; un texte en marge gauche, à demi biffé, s'y rattache : « alors $\mathbb{W}^{a}(f)$ \ill{} \uncertain{morphisme} \uncertain{affine} \ill{} de type fini \ill{} $\mathbb{W}(S)$ \ill{} $\mathbb{W}^{a}(\mathfrak{Y}) \to \mathbb{W}^{a}(\mathfrak{X})$ \ill{} \uncertain{compact} \ill{} ».}

\page{69}
\section{(3) Recollement}

\ill{} \uncertain{hyp.} \uncertain{faites} sur $\mathbb{W}$.

\textbf{Lemme 1} \struck{Proposition} Soit $\mathfrak{X}$ un
$\mathbb{W}(S)$-préschéma, \struck{\ill{}}
\struck{\ill{} \ill{} \ill{} \ill{} $\mathfrak{X}$, \ill{}}
\note{les lignes suivantes sont prises dans un encadré et barrées d'une longue diagonale ; on lit :}
\struck{\uncertain{Supposons} \ill{} \ill{} $T$ sur $S$, et
$\mathbb{W}(T) \xrightarrow{f} \mathfrak{X}$, $T = \bigcup T_i$ (les $T_i$
\ill{}), \ill{} $f(\mathbb{W}(T_i)) \subset \mathfrak{X}_i$. Si les}
\add{$\mathbb{W}(S)$} \ill{} \uncertain{morphisme} \ill{}
\struck{$\mathbb{W}(T)$ \ill{} \uncertain{réunion}}

\ill{} \ill{} \ill{} $\mathfrak{X}$, les \ill{}
$\mathbb{W}^{a}(\mathfrak{U})$ \ill{} \ill{}
$\mathbb{W}^{a}(\mathfrak{X}) = X$ existe, soit $\mathfrak{U}'$ l'\ill{}
\ill{} $\mathfrak{U}'$ \ill{}
\[
  \mathbb{W}(X) \xrightarrow{\ \varphi\ } \mathfrak{X}
  \qquad \mathfrak{U}' \subset \mathfrak{U} .
\]

\textbf{($\alpha$)} \uncertain{Supposons} $\mathbb{W}$ \ill{} : qu'il existe
un \ill{} $\mathfrak{U}$ dans $X$ tel que pour tout $T/S$, \ill{}
\uncertain{morphisme} $\mathbb{W}(T) \xrightarrow{f} \mathbb{W}(X)$, \ill{}
la condition $f(\mathbb{W}(T)) \subset \mathfrak{U}'$ $\Longleftrightarrow$
$f(\mathbb{W}(T)) \subset \mathbb{W}(\mathfrak{U})$ \ill{} [\ill{}
$\mathbb{W}(\mathfrak{U}) \subset \mathfrak{U}'$] \ill{} \ill{}, $\mathfrak{U}$
\ill{} \ill{} $\mathbb{W}(\mathfrak{U}) \to \mathfrak{U}$ \ill{}
\ill{} \ill{} \ill{} $\mathbb{W}(\mathfrak{U}')$.
\note{texte fortement abrégé ; seules les formules sont sûres.}

\page{70}
C'est trivial. \struck{\ill{} \ill{} $\Gamma'$ ou $\Gamma''$}
\note{une note encadrée, dans un cadre en escalier, renvoie à « $\Gamma'$ ou $\Gamma''$ » ; elle est barrée.}

\textbf{\struck{Proposition}} Si $\mathbb{W}^{a}(\mathfrak{X})$ existe et
\ill{}

\textbf{Lemme 2} Soient $\mathfrak{X}/\mathbb{W}(S)$, \ill{} \ill{}
$\mathfrak{X}$ : \ill{} \ill{} \uncertain{ouverts} \ill{}
\struck{\ill{} $\mathbb{W}^{a}(\mathfrak{X}_i)$ \uncertain{existe}}
\struck{$\mathfrak{X}_i$ dans $\mathfrak{X}$ \ill{} \ill{} \ill{}}
\struck{\ill{} \ill{} \ill{}} \ill{} les $\mathbb{W}^{a}(\mathfrak{X}_i)$
\ill{} \uncertain{existent} et \ill{} pour \uncertain{deux}
$\mathfrak{X}_i \cap \mathfrak{X}_j$ \uncertain{ouverts},
$\mathfrak{X}_i$ \ill{} \ill{} \ill{}, [\uncertain{donc}
$\mathbb{W}^{a}(\mathfrak{X}_i \cap \mathfrak{X}_j)$ \uncertain{existe}]
et \ill{} \ill{} \ill{} \ill{} \ill{} \uncertain{(et kif-kif pour les
$\mathfrak{X}_i \cap \mathfrak{X}_j \cap \mathfrak{X}_k$ dans
$\mathfrak{X}_i \cap \mathfrak{X}_j$)} \ill{} \ill{}
$\mathbb{W}^{a}(\mathfrak{X}_i)$ ], \ill{} \ill{}

\textbf{($\beta$)} \uncertain{Pour} \ill{} $T/S$, et \ill{}
\uncertain{morphisme}
\[
  \mathbb{W}(T) \xrightarrow{\ f\ } \mathfrak{X}, \qquad
  T = \bigcup_{i} T_i \ \text{(les \ill{} de $T$) tel que}
\]
$T_i$ \ill{} \ill{} \ill{}
\[
  f(\mathbb{W}(T_i)) \subset \mathfrak{X}_i .
\]
\ill{} \ill{} \uncertain{conditions}, $\mathbb{W}^{a}(\mathfrak{X})$
existe, \ill{} \ill{} \uncertain{réunion} des $\mathbb{W}^{a}(\mathfrak{X}_i)$,
\uncertain{qui} \ill{} \uncertain{identifiés} : \ill{} \ill{} \ill{} à
$\mathbb{W}^{a}(\mathfrak{X})$,

\page{71}
$\mathbb{W}^{a}(\mathfrak{X}_i \cap \mathfrak{X}_j)$ \ill{}
$\gamma_{ij}$ \add{\uncertain{noté}} \ill{} \uncertain{bijection}
\add{\ill{}} \struck{\ill{} $X_{ij}$} \ill{} de
$\mathbb{W}^{a}(\mathfrak{X}_i)$ \ill{} \add{$X_{ji}$ \ill{}}
$\mathbb{W}^{a}(\mathfrak{X}_j)$ : \ill{} \ill{} \ill{}
\uncertain{recollement} \ill{} $\mathbb{W}^{a}(\mathfrak{X}_i) = X_i$. Il
faut vérifier \ill{} \ill{} \struck{$X_{ijk} \ill{} X_{ikj}$} \ill{}
\uncertain{compatibilité} \ill{},
\[
  \begin{array}{ccccc}
  X_{ijk} & \xrightarrow{\varphi_{ji} \mid} & X_{jik} & \xrightarrow{\varphi_{kj} \mid} & X_{kij} \\
  \| & & \| & & \| \\
  X_{ij} \cap X_{ik} & & X_{ji} \cap X_{jk} & & X_{ki} \cap X_{kj}
  \end{array}
\]
\note{une flèche courbe, étiquetée $\varphi_{ki} \mid$, joint au-dessus le premier terme au dernier ; les indices sont lus au plus près.}

Or \uncertain{trivial} \ill{} \uncertain{faire} \ill{} \ill{} \ill{}
\[
  X_{ijk} \simeq \mathbb{W}^{a}(\mathfrak{X}_i \cap \mathfrak{X}_j \cap
  \mathfrak{X}_k)
\]
et \ill{} les applications $\varphi$ \ill{} (\uncertain{p.\ ex.})
s'identifient à l'identité.

On a ainsi défini \ill{} $X = \bigcup X_i$, $X_i \simeq
\mathbb{W}^{a}(\mathfrak{X}_i)$. On définit
\[
  \mathbb{W}(X) \to \mathfrak{X}
\]
en notant que $\mathbb{W}(X) = \bigcup_i \mathbb{W}(X_i)$ \ill{} \ill{}
\ill{} \ill{} \ill{} \uncertain{morphismes}
\[
  \mathbb{W}(X_i) \to \mathfrak{X}_i \subset \mathfrak{X}
\]
\uncertain{par} définition de $X_i = \mathbb{W}^{a}(\mathfrak{X}_i)$. Le

\page{72}
\uncertain{recollement} de ces morphismes \ill{} \uncertain{immédiat}.
Par suite, on a défini $\mathbb{W}(X) \xrightarrow{\ \varphi\ }
\mathfrak{X}$.

Je \uncertain{dis} que c'est une solution du pb universel \ill{} i.e.\ que
pour tout $T$,
\[
  \operatorname{Hom}_{S}(T, X) \to
  \operatorname{Hom}_{\mathbb{W}(S)}(\mathbb{W}(T), \mathfrak{X})
\]
est \uncertain{bijectif}. Soit \struck{$u\colon T \to X$},
\[
  \mathbb{W}(T) \xrightarrow{\ \psi\ } \mathfrak{X}
\]
je dis \uncertain{qu'il} \ill{} \ill{} factorisation de façon unique en
\[
  \mathbb{W}(T) \xrightarrow{\mathbb{W}(u)} \mathbb{W}(X)
  \xrightarrow{\ \varphi\ } \mathfrak{X}
\]
En effet, \struck{\ill{} \ill{} \ill{} \ill{} \ill{} factorisation, \ill{}
\ill{} $T_i$ \ill{}} \struck{$\mathbb{W}(T)$ \ill{}}
utilisons la condition ($\beta$) : $T = \bigcup T_i$, \ill{}
$\psi(\mathbb{W}(T_i)) \subset \mathfrak{X}_i$. Donc
$\psi \mid \mathbb{W}(T_i)$ est défini par $u_i\colon T_i \to X_i$, par
factorisation
\[
  \mathbb{W}(T_i) \xrightarrow{\mathbb{W}(u_i)} \mathbb{W}(X_i)
  \xrightarrow{\ \varphi_i\ } X_i
\]
\note{la dernière flèche aboutit à « $X_i$ » sur la page, où l'on attendrait $\mathfrak{X}_i$.}
On vérifie \ill{} \ill{} que les morphismes $T_i \to X_i \to X$
\struck{$\mathbb{W}(T_i) \to \mathbb{W}(X_i) \to \mathbb{W}(X)$}

\page{73}
\uncertain{coïncident} dans les $T_i \cap T_j$, donc définissent $T \to X$
et \ill{} \ill{} \uncertain{satisfait} \ill{} \uncertain{conditions}
\uncertain{voulues}. D'ailleurs, \struck{\ill{}} \ill{} \ill{} \ill{},
\uncertain{a priori}, \ill{} \ill{} \ill{} construction \ill{} \ill{} \ill{}
\uncertain{fait} \uncertain{précéder}. \uncertain{Introduisons} \ill{}
\ill{} \uncertain{recouvrements} : $(T_i)$, $(T'_j)$, \ill{}
\uncertain{recouvrement} \ill{} $(T_i \cap T'_j)_{(i,j) \in I \times I'}$
\ill{} \ill{} l'\uncertain{unicité}.

\textbf{Lemme 3} La condition ($\beta$) est automatiquement vérifiée
(\uncertain{pour} \ill{} $\mathfrak{X} = \bigcup \mathfrak{X}_i$ \ill{}
\ill{}), \ill{} \ill{} il \ill{} \ill{} un isomorphisme \ill{} de foncteurs
\[
  |T| \xrightarrow{\ \sim\ } |\mathbb{W}(T)|
\]
\uncertain{en} particulier dans les cas II, \struck{\ill{}}, en
particulier III. C'est trivial.

\textbf{Lemme 4} La condition $\alpha$ est automatiquement vérifiée dans
les cas I$'$, II$^{*}$.

C'est trivial pour II$^{*}$. Pour I$'$, \uncertain{prenons}
\uncertain{voisinages} \ill{}, utilisons \ill{} \ill{} un morphisme fini est
\uncertain{fermé}. (\ill{} hyp.\ \uncertain{noeth.})

\page{74}
\struck{\ill{}} \uncertain{Conditions} \ill{} \uncertain{satisfaites} dans le
cas I, si $S'$ est propre sur $S$ et si $\mathbb{W}^{a}(\mathfrak{X})$ est
loc.\ noeth.

\uncertain{Mauvaises} \ill{} \ill{} \ill{} I$'$.

\uncertain{Mettons} ensemble les résultats obtenus

\textbf{Théorème 1} Dans le cas d'un foncteur de Greenberg
\add{\uncertain{spécial}}, ou d'un foncteur $\mathbb{W}_{S'}$ défini par un
morphisme $S' \to S$ fini, loc.\ libre, \struck{\ill{}} \uncertain{radiciel},
\uncertain{le} $\mathbb{W}^{a}(\mathfrak{X})$ existe pour tout préschéma
$\mathfrak{X}$ sur $\mathbb{W}(S)$. \struck{\ill{}} $\mathbb{W}^{a}$ est de
plus un foncteur de \emph{\uncertain{nature} locale}.
\marginal{un petit schéma encadré, illisible \ill{}}

\textbf{Théorème 1 bis} \struck{Dans le cas \ill{}}, I$'$,
\add{\uncertain{Supposons} \ill{} \ill{} dans le cas I$'$.}
\struck{$\varphi_{S'/S}(\ill{})$} (i) Si $\varphi_{S'/S}(\mathfrak{X})$
existe, alors pour tout \ill{} $\mathfrak{U}$ de $\mathfrak{X}$,
$\varphi_{S'/S}(\mathfrak{U})$ existe et $\varphi_{S'/S}(\mathfrak{U}) \to
\varphi_{S'/S}(\mathfrak{X})$ est une immersion \uncertain{ouverte}.

(ii) Soit $\mathfrak{X} = \bigcup \mathfrak{X}_i$, où les $\mathfrak{X}_i$
sont tels que pour tout $s \in S$, tout \struck{\ill{}} \add{\ill{} \ill{}}
de $\mathfrak{X}_s$ sur $S'_s$ soit contenu dans un $\mathfrak{X}_i$.
Alors si les $\varphi_{S'/S}(\mathfrak{X}_i)$ existent,
$\varphi_{S'/S}(\mathfrak{X})$

\page{75}
existe et est \uncertain{réunion} des $\varphi_{S'/S}(\mathfrak{X}_i)$.

\textit{Démonstration.} Il suffit de prouver (ii), et pour cela de prouver
que l'hypothèse faite \uncertain{implique} la condition de recollement. Or,
\uncertain{si} $u\colon T' \to \mathfrak{X}$ est donné, on introduit les
$T'_i = u^{-1}(\mathfrak{X}_i)$, \ill{} et les
\[
  T_i = \complement\, \varphi_{T}(T' - T'_i)
\]
\ill{} \ill{} \ill{} \ill{} de $T$, \ill{} \ill{} prouver que
$T = \bigcup T_i$, \struck{et pour cela, il} \ill{} $t \in T_i$ signifie que
la restriction à $T'_t$ de $\varphi$ \uncertain{applique} $T'_t$ dans
$\mathfrak{X}_i$ : donc \ill{} \ill{} de $(\mathfrak{X}_s/S'_s)_{/k(s)}$
\uncertain{définie} \ill{} l'extension $k(t)$ de $k(s)$, doit être
\uncertain{contenue} dans \ill{} \ill{} $\mathfrak{X}_i$. Or
\struck{si \ill{} \ill{} \ill{} \ill{} $\mathfrak{X}_s$} \add{\ill{} des
$\mathfrak{X}_i^{\alpha}$} \struck{\ill{} \ill{} \ill{} \ill{}
$S'^{\alpha}_s$, \ill{} \ill{} \ill{} \ill{} \ill{} \ill{} \ill{}
$\mathfrak{X}_i^{\alpha}$ \ill{}}
\note{tout ce passage est encadré et barré de diagonales ; au-dessus, un ajout « en des $\mathfrak{X}_i^{\alpha}$ », et à sa droite une insertion « \ill{} $s \in S(\ill{})$ \ill{} ».}

\marginal{diagramme, à gauche des lignes : $\mathfrak{X} \leftarrow T'$
au-dessus de $S \leftarrow S'$ et $T \leftarrow T'$ (flèche $\varphi_T$),
puis $\mathfrak{X}_s$, $S'_s$ sur $k(s)$ et $T'_t$ sur $k(t)$}

\textbf{Corollaire 1} Si \uncertain{toute} partie \struck{\ill{}} finie
\add{\uncertain{d'un}} $\mathfrak{X}_s$ est contenue dans un \ill{}
\uncertain{affine} de $\mathfrak{X}$, \struck{\ill{}} [\uncertain{par}
\uncertain{exemple} si \struck{\ill{}} \add{$\mathfrak{X}/S$} \ill{} \ill{}
\uncertain{quasi-projectifs}] alors $\varphi_{S'/S}(\mathfrak{X})$ existe,
et est
\marginal{\uncertain{Soit} \ill{} \struck{\ill{}} \ill{} \ill{} de $S'/S$}

\page{76}
réunion des $\varphi_{S'/S}(\mathfrak{U})$, où $\mathfrak{U}$
\uncertain{parcourt} les \uncertain{ouverts} affines de $\mathfrak{X}$.

\textbf{Corollaire 2} \uncertain{Supposons} \struck{de plus}
$\mathfrak{X}$ \add{\uncertain{sur} $S'$} de type
\add{\uncertain{tel} que toute partie finie de $\mathfrak{X}$ soit contenue
dans un ouvert affine, et} fini \uncertain{sur} $S$ loc.\ noeth. Alors
$\varphi_{S'/S}(\mathfrak{X})$ est de type fini sur $S$.

La question étant locale \add{sur $S$}, on peut supposer \ill{}
\uncertain{affine}. Tout \uncertain{résulte} \ill{}

\textbf{Lemme} Soit $S$ noethérien, $\mathfrak{X}$ de type fini sur $S$,
\ill{} \ill{} \ill{} \add{\uncertain{les} $(\mathfrak{X}_i)_{i \in I}$
\uncertain{telles} \ill{}} \uncertain{supposons} que \ill{} partie
\struck{finie} $E$ de $\mathfrak{X}$ telle que $\operatorname{card}(E) \le n$
\ill{} contenue dans un \ill{} \struck{\ill{}} $\mathfrak{X}_i$. Alors il
existe une sous-famille finie de $(\mathfrak{X}_i)_{i \in I}$ ayant la
\ill{} propriété.

En effet, l'hyp.\ signifie que
\[
  \bigcup_{i \in I} \underbrace{\mathfrak{X}_i \times_S \cdots \times_S
  \mathfrak{X}_i}_{n\ \text{facteurs}}
  = \underbrace{\mathfrak{X} \times_S \cdots \times_S \mathfrak{X}}_{n\
  \text{facteurs}} .
\]
\ill{} \uncertain{résulte} \ill{} \uncertain{fait} que
$\mathfrak{X} \times_S \cdots \times_S \mathfrak{X}$ est \add{un espace}
\uncertain{noethérien}.

\ill{} Le foncteur $\mathbb{W}^{a}$ \ill{} \ill{} \uncertain{local} en tous
les \ill{} \ill{} \ill{} \ill{} \uncertain{vérifie} \ill{} \ill{}
\note{la dernière ligne est séparée par un trait ; elle est écrite serrée au bord inférieur.}

\page{77}
\section{4) Applications aux immersions, etc.}

\textbf{\struck{Proposition}} Supposons que $\mathbb{W}$ soit
\struck{\ill{}} \add{I$'$, ou Greenberg spécial}.
Soit $f\colon \mathfrak{X} \to \mathfrak{Y}$ un morphisme
\struck{\uncertain{localement} de type fini} d'immersion
\struck{\ill{}} fermée. \add{\uncertain{Supposons} dans le} \ill{}
\uncertain{cas} I$'$ \ill{} \add{\uncertain{l'hypothèse}
\uncertain{faite} \ill{}} \struck{$\mathbb{W}^{a}(\mathfrak{Y})$}
\uncertain{existe} \ill{} \ill{} $\mathfrak{Y}_s$ \ill{}
\struck{\ill{} \ill{} \ill{} $\mathbb{W}^{a}(\mathfrak{U})$ $\mathfrak{U}$}
\add{\uncertain{contenue} dans un \ill{} affine de $\mathfrak{Y}$}.
\struck{\uncertain{affine} \ill{} $\mathfrak{Y}$.} Alors
$\mathbb{W}^{a}(\mathfrak{X})$ existe \add{\uncertain{(existe et)}} et
$\mathbb{W}^{a}(f)$ est une immersion fermée.

En effet, \ill{} \ill{} \ill{} \ill{} \ill{} \uncertain{cas} affine, où
\ill{} \ill{} \ill{}.

\textbf{Corollaire} \add{Sous les conditions précédentes sur $\mathbb{W}$},
si $\mathfrak{Y}$ est \uncertain{quasi-projectif} sur $S^{*}$, et si
$f\colon \mathfrak{X} \to \mathfrak{Y}$ est une immersion, alors
$\mathbb{W}^{a}(f)$ est une immersion.

En effet, $f$ se factorise en
\[
  \mathfrak{X} \xrightarrow{\ f'\ } \mathfrak{X}' \xrightarrow{\ f''\ }
  \mathfrak{Y}
\]
\ill{} $f'$ une immersion fermée et $f''$ immersion
\uncertain{ouverte}. $\mathfrak{X}'$ \ill{} \ill{} \ill{}
\uncertain{quasi-projectif}, \ill{}

\page{78}
\uncertain{appliquer} la \uncertain{prop.} à $f'$, et on \uncertain{applique}
\ill{} à $f''$.

\struck{\ill{}} \textbf{Corollaire 2} Sous les conditions de la prop.\ sur
$\mathbb{W}$, \struck{\ill{}} supposons que \struck{$\mathfrak{X}
\times_{\mathfrak{Y}} \mathfrak{X}$} toute partie finie de $\mathfrak{X}$
soit contenue dans un ouvert affine. Alors si $f\colon \mathfrak{X} \to
\mathfrak{Y}$ est séparé, alors \ill{} \ill{} \ill{} de $\mathbb{W}^{a}(f)$
[\ill{}, \uncertain{supposons} que $\mathbb{W}^{a}(\mathfrak{Y})$ existe dans
le cas I$'$].

On \uncertain{applique} la prop.\ à : $\mathfrak{X} \to \mathfrak{X}
\times_{\mathfrak{Y}} \mathfrak{X}$.

\textbf{Proposition} $\mathbb{W}$, \struck{\ill{}} comme dans la prop.
\uncertain{précédente}. Soit $f\colon \mathfrak{X} \to \mathfrak{Y}$, on
\uncertain{suppose} \add{dans le cas I$'$} \struck{\ill{}} que \ill{}
\uncertain{toute} \ill{} \ill{} \uncertain{ouvert} affine $\mathfrak{U}$ de
$\mathfrak{Y}$, \add{\uncertain{toute} partie finie \ill{}}
$f^{-1}(\mathfrak{U})$ est \struck{\ill{}} contenue dans un \ill{} affine
(\uncertain{p.\ ex.} $f$ \uncertain{quasi-projectif}). On suppose de plus
$\mathfrak{Y}$ \struck{\ill{}} loc.\ noeth.\ et $f$ de type fini. Alors
$\mathbb{W}^{a}(f)$ est de type fini \struck{\ill{} affine}.
\marginal{\uncertain{Dans} le « \uncertain{Greenberg spécial} » \ill{}
\ill{} \ill{} \uncertain{remplacer} \ill{} \ill{} \ill{}
\uncertain{quasi-projectif}}

\page{79}
\textbf{Proposition} Comme dessus, pour un morphisme \emph{affine}.
\note{la page s'arrête sur cet énoncé ; le reste de la feuille est blanc.}

\section{Schémas localement compacts}
\note{titre écrit de sa main, au coin supérieur droit d'une feuille de couverture (p. 80, sans autre contenu), sous un premier titre biffé et illisible \struck{\ill{}} ; il annonce la suite, qui commence au lot suivant.}

\end{document}
